%% Fasit til kortsett fra mattelab.no – laget automatisk av bygg.py.
%% Kompiler med lualatex.
\documentclass[11pt,a4paper]{article}
\usepackage[norsk]{babel}
\usepackage[a4paper,margin=1.2cm,bottom=1.7cm]{geometry}
\usepackage{fontspec}
\usepackage{unicode-math}
\setmainfont{TeX Gyre Pagella}
\setsansfont{TeX Gyre Heros}[Scale=MatchLowercase]
\setmathfont{TeX Gyre Pagella Math}
\usepackage{xcolor}
\definecolor{blekk}{HTML}{16213A}
\definecolor{dempet}{HTML}{566077}
\definecolor{m1}{HTML}{2F4AB0}
\usepackage{tikz}
\usepackage{enumitem}
\usepackage{qrcode}
\usepackage{fancyhdr}
\usepackage[hidelinks]{hyperref}
\pdfvariable suppressoptionalinfo \numexpr 1+2+4+8+16+32+64+128+256+512\relax
\setlength{\parindent}{0pt}
\setlength{\parskip}{0pt}
\pagestyle{fancy}
\fancyhf{}
\renewcommand{\headrulewidth}{0pt}
\fancyfoot[L]{\sffamily\footnotesize\color{dempet}mattelab.no · Fire former av samme funksjon · kortsett K8 · fasit}
\fancyfoot[R]{\sffamily\footnotesize\color{dempet}Side \thepage}
\begin{document}
\begin{minipage}[t]{0.8\linewidth}\vspace{0pt}
{\sffamily\small\color{dempet}Fasit for læreren · kortsett K8 · 1T}\par\vspace{0.25em}
{\fontsize{20}{24}\selectfont\bfseries Fire former av samme funksjon}\par\vspace{0.5em}
Finn fire og fire kort som hører til samme andregradsfunksjon: utvidet form, faktorisert form, nullpunkter og graf.\par\vspace{0.4em}{\small\color{dempet}Kortene er nummerert i den rekkefølgen de står i utskriften.}
\end{minipage}\hfill
\begin{minipage}[t]{0.15\linewidth}\vspace{0pt}\raggedleft
\qrcode[height=1.9cm]{https://mattelab.no/varierende/fire-former.html}
\end{minipage}\par\vspace{0.8em}

\begin{enumerate}[leftmargin=1.6em, itemsep=0.7em]
\item \begin{minipage}[t]{\linewidth}\textsf{\small K8-01}\quad $f(x) = x^2 - x - 6$ \\ \textsf{\small K8-07}\quad \emph{Figur:} Graf til en andregradsfunksjon som åpner oppover. Den skjærer x-aksen i x = −2 og x = 3 og y-aksen i y = −6. \\ \textsf{\small K8-14}\quad Nullpunkter: $x = -2$ og $x = 3$ \\ \textsf{\small K8-16}\quad $f(x) = (x + 2)(x - 3)$\end{minipage}
\item \begin{minipage}[t]{\linewidth}\textsf{\small K8-05}\quad \emph{Figur:} Graf til en andregradsfunksjon som åpner oppover. Den skjærer x-aksen i x = −3 og x = 2 og y-aksen i y = −6. \\ \textsf{\small K8-06}\quad $f(x) = (x + 3)(x - 2)$ \\ \textsf{\small K8-08}\quad Nullpunkter: $x = -3$ og $x = 2$ \\ \textsf{\small K8-21}\quad $f(x) = x^2 + x - 6$\end{minipage}
\item \begin{minipage}[t]{\linewidth}\textsf{\small K8-02}\quad \emph{Figur:} Graf til en andregradsfunksjon som åpner oppover. Den skjærer x-aksen i x = 1 og x = 4 og y-aksen i y = 4. \\ \textsf{\small K8-22}\quad $f(x) = (x - 1)(x - 4)$ \\ \textsf{\small K8-23}\quad $f(x) = x^2 - 5x + 4$ \\ \textsf{\small K8-28}\quad Nullpunkter: $x = 1$ og $x = 4$\end{minipage}
\item \begin{minipage}[t]{\linewidth}\textsf{\small K8-13}\quad $f(x) = x^2 + 5x + 4$ \\ \textsf{\small K8-17}\quad Nullpunkter: $x = -4$ og $x = -1$ \\ \textsf{\small K8-19}\quad \emph{Figur:} Graf til en andregradsfunksjon som åpner oppover. Den skjærer x-aksen i x = −4 og x = −1 og y-aksen i y = 4. \\ \textsf{\small K8-24}\quad $f(x) = (x + 4)(x + 1)$\end{minipage}
\item \begin{minipage}[t]{\linewidth}\textsf{\small K8-11}\quad $f(x) = x^2 - 3x$ \\ \textsf{\small K8-12}\quad Nullpunkter: $x = 0$ og $x = 3$ \\ \textsf{\small K8-26}\quad \emph{Figur:} Graf til en andregradsfunksjon som åpner oppover. Den skjærer x-aksen i x = 0 og x = 3 og y-aksen i y = 0. \\ \textsf{\small K8-27}\quad $f(x) = x(x - 3)$\end{minipage}
\item \begin{minipage}[t]{\linewidth}\textsf{\small K8-03}\quad $f(x) = (x + 2)(x - 2)$ \\ \textsf{\small K8-09}\quad \emph{Figur:} Graf til en andregradsfunksjon som åpner oppover. Den skjærer x-aksen i x = −2 og x = 2 og y-aksen i y = −4. \\ \textsf{\small K8-10}\quad $f(x) = x^2 - 4$ \\ \textsf{\small K8-18}\quad Nullpunkter: $x = -2$ og $x = 2$\end{minipage}
\item \begin{minipage}[t]{\linewidth}\textsf{\small K8-04}\quad Ett nullpunkt: $x = 3$ \\ \textsf{\small K8-15}\quad \emph{Figur:} Graf til en andregradsfunksjon som åpner oppover. Den skjærer x-aksen i x = 3 og y-aksen i y = 9. \\ \textsf{\small K8-20}\quad $f(x) = (x - 3)^2$ \\ \textsf{\small K8-25}\quad $f(x) = x^2 - 6x + 9$\end{minipage}
\end{enumerate}
\subsection*{Til læreren}
Gruppene med nullpunkter −2 og 3 og nullpunkter 2 og −3 ligner hverandre med vilje. Fortegnet i parentesen er motsatt av fortegnet til nullpunktet. Spør også: Hvordan ser du konstantleddet på grafen? Hvorfor har (x − 3)² bare ett nullpunkt?
\end{document}
